\section{The Model}
\label{sec:model}

We consider a target thin region $A$ and an outside region $B$. A local government in region $A$ decides whether to establish a public co-creation hub and, if a hub is in place, whether to supplement it with a cross-regional pipeline policy. The model preserves the paper's two-period structure, but it now distinguishes explicitly between \emph{local coordination within the target region} and \emph{cross-regional renewal from outside the region}. This is the minimal extension needed to evaluate hub policy as a regional policy bundle rather than as a stand-alone local intervention.

\subsection{Environment: Two Regions and Potential Interactions}

Region $A$ contains $N_A \ge 2$ local participants, where a participant may be interpreted as a firm, researcher, entrepreneur, or other actor capable of generating collaborative surplus. Region $B$ contains $N_B \ge 1$ outside participants who form the external pool from which new ideas, collaborators, and partnerships may arrive.

The number of unordered potential \emph{local} pairs within region $A$ is
\begin{equation}
P_A := \frac{N_A(N_A-1)}{2},
\label{eq:pair_count}
\end{equation}
and the number of potential \emph{cross-regional} links between regions $A$ and $B$ is
\begin{equation}
Q_{AB} := N_A N_B.
\label{eq:cross_pair_count}
\end{equation}

A benevolent local planner in region $A$ chooses whether to establish a public hub. Let
\[
h \in \{0,1\}
\]
denote the hub decision, where $h=1$ means that a public hub is established and $h=0$ means that it is not.

If a hub exists, the planner may also implement a cross-regional pipeline policy,
\[
b \ge 0,
\]
which captures efforts to connect region $A$ more intensively to outside participants in region $B$. Examples include guest programs, interregional brokerage, visiting mentors, university--industry linkage schemes, and outreach to firms or researchers located outside the target region.

Time is discrete and consists of two periods, $t=1,2$. The hub affects local matching in both periods. The pipeline policy affects only period 2.

\subsection{Local Coordination Through a Hub}

Let $m_A(h)$ denote the effective local matching rate among potential pairs within region $A$. We assume
\begin{equation}
m_A(h) = m_A^0 + \bigl(m_A^1-m_A^0\bigr)h,
\qquad
0 \le m_A^0 < m_A^1 \le 1.
\label{eq:matching_rate}
\end{equation}
Thus, establishing a hub lowers local interaction costs in region $A$ and raises the fraction of potential local pairs that become realized economically meaningful interactions.

Let $v_L>0$ denote the average social surplus generated by one realized and still-novel \emph{local} interaction. Then period-1 gross surplus is
\begin{equation}
B_1(h) = v_L P_A m_A(h).
\label{eq:B1}
\end{equation}

Equation \eqref{eq:B1} captures the hub's immediate coordination role inside the target region: by lowering local search and meeting frictions, it increases the number of productive local encounters in the short run.

\subsection{Dynamic Redundancy in Local Interactions}

The same local interactions that create value in period 1 may reduce novelty in period 2 if the local network becomes too repetitive. To capture this with minimal structure, let the period-2 \emph{local} component of surplus be
\begin{equation}
B_{2L}(h)
=
\beta v_L P_A m_A(h)\bigl[1-\delta m_A(h)\bigr],
\qquad
\beta\in(0,1],\; \delta>0.
\label{eq:B2_local}
\end{equation}

The parameter $\delta$ measures the rate at which intensive local interaction generates redundancy, repeated exposure, or local lock-in. A higher $\delta$ means that a denser round of local interaction in region $A$ reduces second-period novelty more sharply.

Expression \eqref{eq:B2_local} may be read as a reduced-form summary of repeated local matching. When the hub raises $m_A(h)$, it increases both the number of realized local interactions and the probability that some of those interactions become less novel in period 2 because they draw repeatedly on the same local pool.

We maintain the regularity condition
\begin{equation}
1-\delta m_A^1 > 0,
\label{eq:regularity}
\end{equation}
which ensures that second-period local novelty remains positive even under the hub.

\subsection{Cross-Regional Pipeline Policy}

In period 2, the planner may also create additional links between region $A$ and outside participants in region $B$. Let the effective cross-regional connection rate (or intensity) be
\begin{equation}
q(b)=q_0+\eta b h,
\qquad
q_0\ge 0,
\quad
\eta>0.
\label{eq:q_b}
\end{equation}
The baseline term $q_0$ captures cross-regional access that exists even without new policy. The term $\eta b h$ captures the idea that the pipeline policy raises region $A$'s outside connectivity only when a hub is present to organize, absorb, and channel those external contacts. Accordingly, if $h=0$, then $q(b)=q_0$ for all $b$, so it is without loss to set $b=0$.

Let $v_X>0$ denote the average social surplus generated by one effective \emph{cross-regional} linkage between regions $A$ and $B$. The period-2 \emph{external} component of surplus is then
\begin{equation}
B_{2X}(h,b)
=
\beta v_X Q_{AB} q(b).
\label{eq:B2_external}
\end{equation}

Combining \eqref{eq:B2_local} and \eqref{eq:B2_external}, total period-2 gross surplus is
\begin{equation}
B_2(h,b)
=
\beta\Bigl\{v_L P_A m_A(h)\bigl[1-\delta m_A(h)\bigr] + v_X Q_{AB} q(b)\Bigr\}.
\label{eq:B2}
\end{equation}
Substituting \eqref{eq:q_b} into \eqref{eq:B2} gives
\begin{equation}
B_2(h,b)
=
\beta\Bigl\{v_L P_A m_A(h)\bigl[1-\delta m_A(h)\bigr] + v_X Q_{AB}\bigl(q_0+\eta b h\bigr)\Bigr\}.
\label{eq:B2_expanded}
\end{equation}

This formulation is the key change relative to the single-region baseline. The hub continues to operate through local coordination in region $A$, while the pipeline policy now has a concrete regional interpretation: it raises the rate of cross-regional connections between the target region and an outside pool.

\subsection{Welfare, Timing, and Useful Benchmark Expressions}

Establishing the hub requires a fixed cost $F>0$. The pipeline policy has convex cost
\[
\frac{\kappa}{2}b^2,
\qquad
\kappa>0.
\]
Because the pipeline policy is relevant only when a hub exists, it is convenient to write its cost as $\frac{\kappa}{2}b^2 h$.

Total social welfare is therefore
\begin{equation}
W(h,b)
=
B_1(h)+B_2(h,b)-Fh-\frac{\kappa}{2}b^2 h.
\label{eq:W_general}
\end{equation}
Using \eqref{eq:B1} and \eqref{eq:B2_expanded}, this can be written as
\begin{equation}
W(h,b)
=
v_L P_A m_A(h)
+
\beta\Bigl\{v_L P_A m_A(h)\bigl[1-\delta m_A(h)\bigr] + v_X Q_{AB}\bigl(q_0+\eta b h\bigr)\Bigr\}
-Fh-\frac{\kappa}{2}b^2 h.
\label{eq:W_expanded}
\end{equation}

The planner's problem is to choose $(h,b)$ to maximize \eqref{eq:W_expanded}, taking $(N_A,N_B,m_A^0,m_A^1,v_L,v_X,\delta,q_0,\eta,F,\kappa)$ as given.

The timing is as follows.

\begin{enumerate}[leftmargin=2em]
    \item At the beginning of period 1, the planner chooses whether to establish a public hub, $h\in\{0,1\}$.
    \item Period-1 local matching occurs in region $A$, and period-1 surplus $B_1(h)$ is realized.
    \item At the beginning of period 2, if $h=1$, the planner chooses the cross-regional pipeline policy $b\ge 0$.
    \item Period-2 local and cross-regional surplus, summarized by $B_2(h,b)$, is realized.
\end{enumerate}

Since there is no uncertainty, this staged timing is equivalent to choosing $(h,b)$ at the outset under commitment. The sequential formulation is adopted only because it highlights the interpretation of $b$ as a policy response to the hub's dynamic tendency toward local redundancy.

For later use, it is convenient to record the welfare levels associated with the three benchmark cases.

First, if no hub is established, then $h=0$ and $b=0$, so
\begin{equation}
W(0,0)
=
v_L P_A m_A^0
+
\beta\Bigl\{v_L P_A m_A^0\bigl(1-\delta m_A^0\bigr) + v_X Q_{AB} q_0\Bigr\}.
\label{eq:W00}
\end{equation}

Second, if a hub is established but no pipeline policy is used, then
\begin{equation}
W(1,0)
=
v_L P_A m_A^1
+
\beta\Bigl\{v_L P_A m_A^1\bigl(1-\delta m_A^1\bigr) + v_X Q_{AB} q_0\Bigr\}
-F.
\label{eq:W10}
\end{equation}

Third, if a hub is established and pipeline policy $b$ is used, then
\begin{equation}
W(1,b)
=
W(1,0)
+
\beta\eta v_X Q_{AB} b
-
\frac{\kappa}{2}b^2.
\label{eq:W1b}
\end{equation}

Subtracting \eqref{eq:W00} from \eqref{eq:W10}, the net welfare gain from a \emph{hub alone} is
\begin{equation}
\Delta_H
:=
W(1,0)-W(0,0)
=
v_L P_A \bigl(m_A^1-m_A^0\bigr)\Bigl[(1+\beta)-\beta\delta\bigl(m_A^1+m_A^0\bigr)\Bigr]-F.
\label{eq:DeltaH}
\end{equation}

Subtracting \eqref{eq:W00} from \eqref{eq:W1b}, the net welfare gain from a \emph{hub with pipeline policy} is
\begin{equation}
\Delta_{H+b}(b)
:=
W(1,b)-W(0,0)
=
\Delta_H
+
\beta\eta v_X Q_{AB} b
-
\frac{\kappa}{2}b^2.
\label{eq:DeltaHb}
\end{equation}

Equations \eqref{eq:DeltaH} and \eqref{eq:DeltaHb} already show the model's basic logic. The stand-alone value of the hub is governed by \emph{local} coordination inside region $A$ and the associated risk of dynamic redundancy. By contrast, the incremental value of the pipeline policy is governed by the strength of \emph{cross-regional} renewal through additional $A$--$B$ linkages. The baseline outside connection rate $q_0$ affects levels of welfare, but it cancels out of the incremental hub-alone comparison.

The next section uses \eqref{eq:DeltaH} to characterize when a public hub is socially efficient as a stand-alone policy. Section 5 then distinguishes short-run success from dynamic efficiency. Section 6 introduces optimal pipeline policy and characterizes when a hub alone is undesirable but a hub combined with cross-regional renewal is socially efficient.
